\documentclass[10pt]{article}

\usepackage[margin=2.3cm]{geometry}
\usepackage[T1]{fontenc}
\usepackage[utf8]{inputenc}
\usepackage{lmodern}
\usepackage{parskip}
\usepackage{enumitem}
\usepackage{hyperref}
\usepackage{xcolor}
\usepackage{titlesec}
\usepackage{longtable}
\usepackage{booktabs}
\usepackage{array}
\usepackage{fancyhdr}
\usepackage{lastpage}

\definecolor{linkblue}{HTML}{1155CC}
\hypersetup{
  colorlinks=true,
  linkcolor=linkblue,
  urlcolor=linkblue,
  citecolor=linkblue,
  pdftitle={NMMA Codebase Review -- Retreat Prep},
  pdfauthor={Claude Code}
}

\titleformat{\section}{\normalfont\Large\bfseries}{\thesection}{0.6em}{}
\titleformat{\subsection}{\normalfont\large\bfseries}{\thesubsection}{0.6em}{}
\titleformat{\subsubsection}{\normalfont\normalsize\bfseries}{\thesubsubsection}{0.6em}{}
\titlespacing*{\section}{0pt}{1.3em}{0.5em}
\titlespacing*{\subsection}{0pt}{1.0em}{0.35em}
\titlespacing*{\subsubsection}{0pt}{0.8em}{0.3em}

\newcommand{\code}[1]{\texttt{#1}}

\pagestyle{fancy}
\fancyhf{}
\fancyhead[L]{NMMA Codebase Review}
\fancyhead[R]{\thepage\ / \pageref{LastPage}}
\renewcommand{\headrulewidth}{0.4pt}
\setlength{\headheight}{14pt}

\setlist[itemize]{topsep=2pt, itemsep=3pt, parsep=0pt}
\setlist[enumerate]{topsep=2pt, itemsep=4pt, parsep=0pt}

\title{\vspace{-2em}NMMA Codebase Review --- Retreat Prep}
\author{Generated by Claude Code}
\date{2026-08-23 \quad \textit{Research only --- no code changed, nothing pushed/merged}}

\begin{document}
\maketitle
\thispagestyle{fancy}

\noindent
\textbf{Scope:} \code{nmma/} source tree, tests, docs, CI/lint config, and the live
GitHub issue/PR tracker for \code{nuclear-multimessenger-astronomy/nmma}. Goal: give
the team a prioritized, evidence-based punch list for a 1-week retreat, split by who
should do the work. This report merges two independent passes over the codebase and
issue tracker, so most findings were cross-checked twice; a couple of counts (e.g.
the number of \code{FIXME} comments) differ slightly between the two passes
depending on exactly what was grepped --- noted inline where relevant.

\textbf{Caveat on one audit pass:} part of this review ran in a Python venv
belonging to a \emph{different} project, so \code{healpy} and \code{flake8} weren't
importable there and a live \code{flake8 --max-complexity 15} pass couldn't be run.
All findings below come from static grep/read/\code{git log} inspection instead ---
solid for locating problems, but a real \code{flake8} run in the actual \code{nmma}
dev env (\code{pip install -e ".[dev,grb,neuralnet]"}) should be the first thing
done at the retreat to corroborate and extend \S4.2.

\section{Headline findings}

\begin{enumerate}
\item \textbf{Five of eight top-level modules have essentially zero dedicated
  tests}: \code{eos/} (1{,}775 lines), \code{gw/} (311 lines), \code{mlmodel/}
  (1{,}333 lines), \code{population/} (28 lines), \code{post\_processing/} (1{,}425
  lines, only touched incidentally via \code{maximum\_mass.py}). That's roughly
  \textbf{4{,}850 lines of core physics/inference code with no correctness test}.
  \code{nmma/core/mpi\_setup.py} (683 lines, the MPI-parallel \code{pbilby\_sampling}
  path used whenever \code{USE\_MPI and sampler == 'dynesty'}) also has zero
  coverage --- exactly the kind of ordering-sensitive distributed code that
  regresses silently. The maintainers already know part of this --- issue
  \textbf{\#196 ``Tests for nmma/eos''} and \textbf{\#299 ``Joint analysis unit
  test''} have been open since \textbf{2023-08} and \textbf{2024-01} respectively.

\item \textbf{\code{nmma/tests/models.py} is entirely \code{pytest.mark.skip}'d}
  (GitLab SVD download path retired) and never replaced with an equivalent test ---
  a whole test file is dead weight that silently reports 0 real coverage while
  looking like it exists. Separately, \textbf{test collection was broken for 4 of
  14 test modules} in one audit environment (\code{maximum\_mass.py},
  \code{systematics.py}, \code{tools.py}, \code{training.py} all failed with
  \code{ModuleNotFoundError: No module named 'healpy'}) --- this traces back to
  \code{nmma/em/\_\_init\_\_.py} eagerly importing the entire \code{em} submodule
  tree, so even a small, unrelated import drags in \code{healpy} and other
  heavy/optional dependencies. Worth re-verifying in a correctly provisioned env,
  but the coupling itself is real and worth loosening regardless.

\item \textbf{Two concrete, currently-shipping bugs, both cheap to fix and both
  untested:}
  \begin{itemize}
  \item \code{nmma/core/base.py}'s \code{constraints} setter ($\sim$lines 55--65):
    if \code{value} isn't a \code{PriorDict}, \code{Constraint}, or \code{dict},
    the local \code{constr} variable is never assigned, and
    \code{self.\_constraints = constr} raises an opaque \code{UnboundLocalError}
    instead of a clear validation error.
  \item \code{nmma/population/pop\_likelihood.py}'s
    \code{NeutronStarPopulation.\_\_init\_\_} only handles
    \code{model\_name in \{'flat', 'peak'\}}; any other value silently leaves
    \code{self.distribution} unset, surfacing later as a confusing
    \code{AttributeError} in \code{log\_likelihood()} rather than a clear error at
    construction time.
  \end{itemize}

\item \textbf{$\sim$50--65 \code{FIXME}/\code{TODO} comments} (counts differ
  slightly by grep pattern across the two audit passes), many self-flagged as
  broken or approximate: \code{nmma/em/lightcurve\_generation.py:145} questions its
  own redshift-correction exponent; \code{nmma/core/conversion.py:549} flags a
  physics inconsistency between BNS and NSBH GRB channels (``NSBH might produce a
  GRB too, why not provide the same expression for BNS?''); \code{nmma/em/io.py:545}
  is tagged ``Legacy??? seems unused.'' Densest clusters:
  \code{nmma/joint/injection\_handling.py} (11, including a self-flagged fragile
  dependency on undocumented bilby internals at line 300) and \code{nmma/em/model.py}
  (10, including \emph{``This is an unclean default, should be more explicit!''} and
  \emph{``This is incomplete, but identical to handling in NMMA 0.2.2''} --- i.e.
  known-incomplete logic shipped on purpose to match old behavior). These are
  landmines: they read as ``someone already knows this is wrong,'' but nobody has
  turned them into tracked issues or tests.

\item \textbf{11 bare \code{except:} clauses} across \code{core/}, \code{em/},
  \code{eos/}, \code{post\_processing/}: \code{nmma/eos/eos\_likelihood.py:137,154},
  \code{nmma/eos/eos\_processing.py:57}, \code{nmma/core/base.py:246,379},
  \code{nmma/core/utils.py:21}, \code{nmma/core/mpi\_setup.py:475},
  \code{nmma/em/training.py:532,584},
  \code{nmma/post\_processing/maximum\_mass\_constraint.py:159},
  \code{nmma/post\_processing/hubble\_estimates.py:137} --- these swallow
  \code{KeyboardInterrupt}/\code{SystemExit} and make debugging silent failures in
  long sampler runs much harder. Related: \textbf{106 \code{print()} calls in
  production code}, inconsistently mixed with real logging (e.g.
  \code{nmma/core/base.py} uses \code{logger.info} in one function and plain
  \code{print()} in \code{bilby\_sampling} at lines 297--298, 308) --- no
  consistent logging discipline.

\item \textbf{\code{nmma/mlmodel/} is the most stale, least tested subsystem in the
  repo.} Untouched since 2024-07-18 (\code{inference.py} since 2025-01-09) while
  the rest of core/em/joint/eos were touched as recently as June--July 2026; zero
  tests; ships 9.6 MB of binary model weights
  (\code{frozen-flow-weights.pth} 8.0 MB, \code{similarity\_embedding\_weights.pth}
  1.6 MB) committed directly to git --- not via LFS, not covered by any
  \code{.gitattributes} binary rule --- added in 2024 and never revisited. Needs an
  explicit decision: maintain it (add tests, document it) or mark it clearly
  experimental/deprecated.

\item \textbf{Docs have two actively-broken pages, plus thin docstrings.}
  \code{doc/index.rst} (the Sphinx ``Installation'' page) and
  \code{doc/lfi\_analysis.md:20} tell new users to
  \code{pip install -r requirements.txt} / \code{pip install -r ml\_requirments.txt}
  --- \textbf{neither file exists in the repo anymore}; the project moved to
  \code{pyproject.toml} extras (\code{pip install -e ".[dev,grb,neuralnet]"}, per
  CLAUDE.md) and the docs were never updated. Separately,
  \code{doc/quick-start-guide.rst} (untouched since 2023-09-19) documents CLI flags
  \code{--aligned-spin} and \code{--injection-outfile} that \textbf{do not exist
  anywhere in the source tree} --- copy-pasting its example command fails with an
  argparse error today. A brand-new user following either page verbatim hits an
  error on step 1. Beyond that: \code{nmma/core/base.py} (436 lines, the shared
  likelihood base every domain extends) has $\sim$4 docstrings;
  \code{nmma/joint/joint\_likelihood.py} has $\sim$3; \code{nmma/em/model.py} (1720
  lines, 30+ model classes) is better but still partial ($\sim$27 docstrings).
  \code{tutorials/} has only 2 notebooks against 25+ documented models and
  $\sim$17 CLI entry points. \code{README.md} (166 lines) has no installation
  instructions, CLI quick-start, or code examples at all --- citation/funding
  content only. \code{doc/models.md} links to a stale line-anchor in
  \code{nmma/em/model.py} (points at line 72; actual content is now $\sim$97--125)
  --- model names themselves are still accurate, only the link target drifted.
  Further drift: issues \textbf{\#374}, \textbf{\#307}, \textbf{\#213}/\textbf{\#136},
  and doc pages tied to the still-open \textbf{\#388} design question.

\item \textbf{Tooling/CI pins are stale and inconsistent.}
  \code{.pre-commit-config.yaml} pins \code{black==22.3.0} (current stable is
  24.x/25.x) and \code{pre-commit-hooks==v3.3.0} (current is v5.x). \code{.flake8}
  globally ignores \code{E501} with \code{max-line-length = 80} set but never
  enforced --- decorative config. \code{.github/workflows/*.yml} mixes
  \code{actions/checkout@v2}/\code{@v3}/\code{@v4}, and
  \code{publish-to-pypi.yml:13} uses \code{actions/checkout@main} (tracks a moving
  branch --- a supply-chain risk specifically in the \emph{publish} workflow), plus
  \code{actions/setup-python@v3} (current is v5).

\item \textbf{PR queue has two multi-megabyte-diff PRs} (\#105, \#82: +3.3M lines
  each) sitting open since June 2023 --- almost certainly stale branches with
  binary/pickle files being diffed as text (no \code{.gitattributes} rule stops
  this). A \code{.gitattributes} rule (e.g. \code{*.pkl -diff -text},
  \code{*.pth -diff -text}) would prevent recurrence.

\item \textbf{Trivial repo hygiene:} \code{doc/\_static/nmma-docs.css\textasciitilde}
  and \code{doc/\_static/.nmma-docs.css.un\textasciitilde} are vim backup/swap
  files checked into git since 2022. \code{nmma/core/utils.py:129} emits a live
  \code{DeprecationWarning} from an un-prefixed regex string
  (\code{sep='\textbackslash s+'} should be \code{sep=r'\textbackslash s+'}).
  \code{loadEventSpec()} in \code{nmma/em/io.py:546} is confirmed dead code
  (already self-flagged, zero references elsewhere).

\item \textbf{Design-principle read: the core abstraction is genuinely sound; the
  joint pipeline's glue code is the weak point.}
  \code{NMMALikelihoodMixin}/\code{NMMALikelihood} (\code{nmma/core/base.py:37-185})
  cleanly define the
  \code{parameter\_conversion}/\code{sub\_log\_likelihood}/\code{sanity\_checks}/
  \code{constraints} contract, extended uniformly by
  \code{GravitationalWaveTransientLikelihood}, \code{EquationofStateLikelihood},
  \code{EMTransientLikelihood} --- not a leaky design. But
  \code{MultiMessengerLikelihood.setup\_from\_args}
  (\code{nmma/joint/joint\_likelihood.py:89-177}) is a 90-line procedural
  dispatcher that self-documents its own fragility inline
  (\code{\# FIXME: Find a better way to check this automatically}). And
  configuration handling mixes four serialization mechanisms end-to-end (CLI
  parsers, pickled data dumps between \code{nmma-generation}/\code{nmma-analysis}
  with \textbf{no schema/version check}, YAML, HDF5) --- the unversioned pickle
  dump is the most fragile as an inter-process contract: a stale dump from before
  a code change can silently misbehave.
\end{enumerate}

\section{Prioritized action items (retreat punch list)}

\subsection*{P0 --- do first (small effort, high leverage, unblocks everything else)}
\begin{enumerate}
\item Stand up \textbf{coverage-gated CI} for \code{eos/}, \code{gw/},
  \code{mlmodel/}, \code{population/}, \code{post\_processing/}, and
  \code{core/mpi\_setup.py} --- even one smoke test per module raises the floor.
  (Ties to \#196, \#299.)
\item Delete or replace \code{nmma/tests/models.py} --- a fully-skipped test file
  is worse than no file.
\item Fix the two concrete latent bugs: the \code{base.py} \code{constraints}
  setter (\code{UnboundLocalError} $\to$ clear \code{ValueError}) and
  \code{NeutronStarPopulation.\_\_init\_\_} (guard unknown \code{model\_name}).
  Both are small, mechanical, and easy to unit-test alongside the fix.
\item Replace the 11 bare \code{except:} clauses $\to$ catch specific exceptions or
  at minimum \code{except Exception:} with a logged message.
\item Fix both broken doc entry points: \code{doc/index.rst}'s Installation
  section and \code{doc/lfi\_analysis.md:20} (replace with
  \code{pip install -e ".[dev,grb,neuralnet]"}), and
  \code{doc/quick-start-guide.rst}'s example command (bad flags $\to$ real,
  working flags). These are the very first things a new user hits --- highest
  leverage-per-minute fix in the whole review.
\item Bump \code{pre-commit} tool pins (\code{black}, \code{pre-commit-hooks}) and
  pin GitHub Actions versions consistently; replace \code{actions/checkout@main}
  in \code{publish-to-pypi.yml} with a pinned version --- a moving-branch checkout
  in a \emph{publish} workflow is a real supply-chain risk.
\item Repo hygiene sweep: remove stray editor files, delete confirmed-dead
  \code{loadEventSpec()}, fix the un-prefixed regex in \code{core/utils.py:129},
  add a \code{.gitattributes} binary-diff rule to stop future mega-diff PRs.
\end{enumerate}

\subsection*{P1 --- do this week (medium effort)}
\begin{enumerate}
\setcounter{enumi}{7}
\item Triage and resolve the $\sim$50--65 \code{FIXME}/\code{TODO} comments,
  prioritizing the two densest clusters (\code{joint/injection\_handling.py},
  \code{em/model.py}).
\item Sweep the 106 stray \code{print()} calls into existing loggers.
\item Close or rebase the two mega-diff stale PRs (\#105, \#82).
\item Resolve version-key duplication: \code{distance} vs
  \code{luminosity\_distance} (\#375).
\item Decide and document a \textbf{versioning/release policy} (\#183, \#429,
  \#184) --- blocked partly by \#335 and \#382.
\item Reconcile CLI docs against \code{[project.scripts]} (18 entry points, 7
  never invoked in any test); fix \code{doc/models.md}'s stale line anchor.
\item Decide the long-term status of \code{nmma/mlmodel/} (maintain-with-tests vs.
  explicitly deprecate).
\end{enumerate}

\subsection*{P2 --- track for next cycle (larger, needs design time)}
\begin{enumerate}
\setcounter{enumi}{14}
\item Resolve \code{core\_model\_name} design ambiguity (\#388).
\item Decouple training/analysis code from the pinned TensorFlow version ceiling
  (\#335, \#334), and address poor training results on the new grid (\#301).
\item Redesign \code{gwem\_resampling}'s KDE/approximation scheme (\#78).
\item Add a schema/version check to the \code{*\_dump.pickle} inter-process
  contract between \code{nmma-generation} and \code{nmma-analysis}.
\item Consider adding a \code{conda}-first \code{environment.yml} (\#206).
\item Investigate the long-standing dynesty/ultranest indefinite-hang bug
  (\#202).
\end{enumerate}

\section{Split by who should do it}

\subsection{Needs human/expert judgment (physics, API design, deprecation calls)}
These require domain expertise, a design decision, or the authority to
break/deprecate an API --- not appropriate to hand to an AI agent unsupervised.

\begin{itemize}
\item Resolve the physics \code{FIXME}s in \code{nmma/core/conversion.py} (BNS vs
  NSBH GRB ejecta fitting, prompt-collapse thresholds,
  \code{dynamic\_mass\_fitting\_KrFo} routing) --- scientific correctness
  questions, not code bugs.
\item \code{nmma/em/lightcurve\_generation.py:145} --- whether the $(1+z)$ vs
  $(1+z)^2$ bolometric luminosity correction is right.
\item Decide the fate of \code{core\_model\_name} (\#388).
\item Decide the long-term status of \code{nmma/mlmodel/} (maintain vs.
  deprecate).
\item Redesign or explicitly accept
  \code{MultiMessengerLikelihood.setup\_from\_args} as the ad-hoc dispatcher it
  currently is.
\item Decide on a schema/versioning strategy for the \code{*\_dump.pickle}
  inter-process contract.
\item Versioning/release policy (\#183, \#429, \#184) and what ``NMMA 1.1'' should
  contain.
\item TensorFlow version ceiling (\#335) and training-quality regression (\#301,
  \#302) --- require actually re-validating training code, a nontrivial
  numerical-behavior check.
\item Redesign of \code{gwem\_resampling}'s KDE scheme (\#78) and the sequential
  Bayesian Hubble-constant update (\#100) --- statistical methodology decisions.
\item Deciding what to do with the two mega-diff stale PRs (\#105, \#82) and the
  2023-era draft PR (\#215).
\item Deciding which of the $\sim$50--65 \code{FIXME}s represent real physics bugs
  vs. acceptable approximations.
\item Coverage strategy for \code{eos/}/\code{gw/}/\code{mlmodel/}/
  \code{mpi\_setup.py} --- what counts as a meaningful test needs a domain expert.
\item Population module correctness (\code{NeutronStarPopulation} only supports 2
  hardcoded distributions) --- generalizing is a design call.
\item The \code{richpool} vs \code{schwimmbad} question (\#432, PR \#433) ---
  touches the untested \code{mpi\_setup.py}; needs MPI validation expertise.
\item The long-standing dynesty/ultranest indefinite-hang bug (\#202) --- deep
  sampler-specific debugging.
\item Milky-Way extinction law (G23) + SFD recalibration (\#426) --- a
  physics-correctness change to shipped results, needs expert sign-off.
\item Reviewing/merging the three large, physics-adjacent open PRs: \#409 (RA/Dec
  units fix, correctness-critical), \#397 (tensor-loading refactor), \#379 (new
  LANL1D model, now with merge conflicts, needs science review).
\end{itemize}

\subsection{AI-agent-tractable (mechanical, well-specified, easy to review)}
These are good candidates to hand to an AI coding agent, with human review of the
(small) diff.

\begin{itemize}
\item Fix both broken doc entry points (install instructions; quick-start-guide's
  bad flags) --- purely mechanical, verifiable by running the corrected commands.
\item Add the missing \code{ValueError} in
  \code{NeutronStarPopulation.\_\_init\_\_}; fix the \code{UnboundLocalError} in
  \code{base.py}'s \code{constraints} setter.
\item Replace the 11 bare \code{except:} with \code{except Exception:} + logging.
\item Sweep the 106 stray \code{print()} calls into existing loggers.
\item Delete/replace the fully-skipped \code{nmma/tests/models.py}.
\item Remove stray editor files; delete confirmed-dead \code{loadEventSpec()}; fix
  the un-prefixed regex in \code{core/utils.py:129}; add a \code{.gitattributes}
  binary-diff rule.
\item Bump \code{black}/\code{pre-commit-hooks} pins and pin GitHub Actions
  versions consistently.
\item Write smoke tests for the untested pure-ish functions in
  \code{post\_processing/}, \code{eos/eos\_processing.py}, and
  \code{core/mpi\_setup.py}. Directly matches \textbf{\#196}/\textbf{\#299} ---
  good agent-assignable tasks once a human specifies expected test cases.
\item Fix arm64 install docs per the working recipe already given in
  \textbf{\#374}.
\item \textbf{\#375} deprecated-key cleanup --- the issue body already lists the
  exact files; final ``which key wins'' call is P1 human sign-off.
\item \textbf{\#382} \code{Bu2022mv} model not found --- narrow, reproducible bug
  report with a code snippet already in the issue.
\item \textbf{\#357} timeshift dropped by \code{processData()} --- bounded bug,
  clear failure mode.
\item \textbf{\#263} DOI-fetch 429 errors --- cache the Zenodo DOI lookup once
  instead of hitting the API from every MPI rank.
\item \textbf{\#81} ``sample times double call'' --- remove the redundant
  \code{sample\_times} argument; redo small rather than resurrect the abandoned
  PR \#82.
\item Small doc cross-reference sweep: confirm every \code{[project.scripts]}
  entry point has a doc mention; fix \code{doc/models.md}'s stale line anchor.
\item Mechanical first pass at missing docstrings in \code{nmma/core/base.py} and
  \code{nmma/joint/joint\_likelihood.py}.
\item Additional narrowly-scoped issues with clear specs: \textbf{\#206} conda
  \code{environment.yml}, \textbf{\#346} JSON lightcurves in example files,
  \textbf{\#307} joint-analysis docs, \textbf{\#136} Zenodo quickstart docs,
  \textbf{\#181} compatibility list, \textbf{\#188} \code{--debug} CLI flag,
  \textbf{\#142} plot multiple models (\code{good first issue}, reference
  implementation linked), \textbf{\#184} release-notes writing (once policy is
  set).
\item PR \textbf{\#276} (apptainer definition file) --- small, self-contained,
  already sanity-checked; just needs a rebase/re-verify and merge.
\end{itemize}

\section{Detailed audit}

\subsection{Test coverage gaps}
\begin{itemize}
\item \textbf{\code{nmma/gw/} has no dedicated test file.} \code{gw\_likelihood.py}
  (247 lines), \code{gw\_inputs.py} (35 lines --- its own header calls it ``a hacky
  subclass to adapt bilby\_pipe data generation''), and \code{gw\_parsing.py} are
  exercised only indirectly through \code{nmma/tests/joint\_analysis\_pipeline.py}.
\item \textbf{\code{nmma/eos/} has no dedicated test file.}
  \code{eos\_likelihood.py} (730 lines, $\sim$12 classes) is only touched
  incidentally via joint-pipeline imports. \code{eos\_gen.py} (408 lines),
  \code{eos\_processing.py} (454 lines), \code{tov.py} (109 lines) have zero test
  references.
\item \textbf{\code{nmma/mlmodel/}} (1{,}333 lines across 5 modules + 2 binary
  weight files) has zero tests and isn't imported by any test file.
\item \textbf{\code{nmma/population/}} has zero tests despite being wired directly
  into \code{MultiMessengerLikelihood.setup\_from\_args} and reachable via
  \code{--population-model} --- and it has a real, untested bug (\S1.3).
\item \textbf{\code{nmma/core/mpi\_setup.py}} (683 lines) has no test coverage at
  all.
\item \textbf{\code{nmma/core/conversion.py}} (1{,}018 lines) is only exercised
  indirectly via \code{nmma/tests/injections.py} --- no direct unit tests of
  \code{cosmology\_to\_distance}, \code{nsbh\_parameter\_conversion},
  \code{bns\_ejecta\_conversion}, \code{grb\_energy\_conversion},
  \code{remnant\_disk\_mass\_fitting}.
\item \textbf{CLI entry points never invoked in any test}:
  \code{lightcurve-injection-slurm-setup}, \code{create-lightcurve-slurm},
  \code{svdmodel-download}, \code{multi-config-analysis},
  \code{plot-svdmodel-benchmarks}, \code{gwem-Hubble-estimate},
  \code{gwem-resampling}.
\item Coverage is heavily skewed toward \code{injections.py} (693 test lines) and
  \code{systematics.py} (24 test functions), while \code{eos} (1{,}775 lines) and
  \code{gw} (311 lines) get essentially none and \code{mlmodel} (1{,}333 lines)
  gets none at all.
\end{itemize}

\subsection{Code quality / design smells}
\begin{itemize}
\item \textbf{11 bare \code{except:} clauses} --- full list in \S1.5.
\item \textbf{106 \code{print()} calls in production code}, inconsistently mixed
  with real logging.
\item \textbf{Large, multi-responsibility files:} \code{em/model.py} (1{,}720
  lines), \code{em/utils.py} (1{,}244 lines), \code{core/conversion.py} (1{,}018
  lines), \code{eos/eos\_likelihood.py} (730 lines). A live \code{flake8
  --max-complexity 15} run would very likely flag these four.
\item \textbf{Binary model weights committed directly to git}, not via LFS:
  \code{frozen-flow-weights.pth} (8.0 MB), \code{similarity\_embedding\_weights.pth}
  (1.6 MB).
\item No Python-2-isms found on spot check --- staleness is concentrated in
  \code{mlmodel}/\code{population}/docs/CI config, not core language usage.
\end{itemize}

\subsection{Staleness}
\begin{itemize}
\item \textbf{\code{nmma/mlmodel/*}} last touched 2024-07-18
  (\code{inference.py} 2025-01-09) --- over a year stale.
\item \textbf{\code{nmma/population/pop\_likelihood.py}} last touched 2026-05-04 ---
  noticeably older than sibling files, consistent with being under-maintained.
\item \textbf{\code{nmma/em/systematics.py}} last touched 2026-02-03, months behind
  \code{em/model.py} (2026-07-04) and \code{em/training.py} (2026-06-30).
\item CI action-version inconsistency and stale pre-commit pins --- see \S1.8.
\end{itemize}

\subsection{Design-principle concerns}
\begin{itemize}
\item \textbf{Positive finding, high confidence:} the core abstraction is
  genuinely well-factored (\S1.11) --- \code{NMMALikelihoodMixin}/
  \code{NMMALikelihood} cleanly define a contract that all domain likelihoods
  extend uniformly; \code{MultiMessengerLikelihood} composes them without
  messenger-specific internals beyond routing by \code{isinstance}. Not a leaky
  design.
\item \code{MultiMessengerLikelihood.setup\_from\_args} is the single most
  important function in the joint pipeline and it reads as ad hoc (\S1.11).
\item Four different serialization mechanisms are used end-to-end for
  configuration (CLI parsers, unversioned pickle dumps, YAML, HDF5) --- the
  pickle dump is the most fragile as an inter-process contract.
\item \code{nmma/em/\_\_init\_\_.py} eagerly imports the entire submodule tree,
  pulling in \code{healpy} and other heavy dependencies for any small import ---
  avoidable coupling, and the direct cause of the test-collection failures in
  \S1.2.
\end{itemize}

\section{GitHub issue triage (36 open, as of 2026-08-23)}

\begingroup
\small
\begin{longtable}{@{}p{0.9cm}p{5.1cm}p{1.2cm}p{2.3cm}p{5.6cm}@{}}
\toprule
\textbf{\#} & \textbf{Title} & \textbf{Age} & \textbf{Class.} & \textbf{Rationale} \\
\midrule
\endhead
432 & Using richpool & 2026-08 & Needs human & Dependency-swap proposal for untested \code{mpi\_setup.py}; paired with PR \#433. \\[3pt]
429 & NMMA 1.1 release? & 2026-07 & Needs human & Release/admin-rights process decision. \\[3pt]
426 & Milky-Way extinction law (G23) + SFD recalibration & 2026-07 & Human (AI can draft) & Physics-correctness change to shipped results. \\[3pt]
400 & \code{default\_prior} in \code{InjectionCreator} & 2025-08 & Needs human & Blocked on upstream bilby/bilby\_pipe design decision. \\[3pt]
388 & Remove reliance on \code{core\_model\_name} & 2024-11 & Needs human & API/design change to SVD model naming/loading. \\[3pt]
382 & \code{Bu2022mv} not found in models list & 2024-08 & AI-tractable & Clear repro snippet already in the issue. \\[3pt]
375 & Change deprecated key \code{distance}$\to$\code{luminosity\_distance} & 2024-08 & AI-tractable & Issue enumerates every file to change. \\[3pt]
374 & Update docs for arm64 & 2024-08 & AI-tractable & Fix already given in the issue. \\[3pt]
357 & Timeshift dropped by \code{processData()} & 2024-04 & AI-tractable & Bounded bug, clear failure mode. \\[3pt]
346 & Add JSON lightcurves to example files & 2024-03 & AI-tractable & Simple additive task. \\[3pt]
335 & Relax tensorflow version requirement & 2024-03 & Needs human & Requires re-validating training against newer TF. \\[3pt]
307 & Joint analysis documentation & 2024-01 & AI-tractable & Pure doc-writing task. \\[3pt]
302 & Auto-train tensorflow models on updated grids & 2024-01 & Needs human & New automation/infra feature. \\[3pt]
301 & Poor tensorflow training results for new grid & 2024-01 & Needs human & ML/training-quality diagnosis. \\[3pt]
299 & Joint analysis unit test & 2024-01 & AI-tractable & Matches a test-coverage gap in this audit. \\[3pt]
263 & DOI not retrievable with too many processes & 2023-10 & AI-tractable & Concurrency bug, well-understood fix. \\[3pt]
257 & Calculate chi2 between observation and grid & 2023-10 & Needs human & New feature, needs design decision on statistic. \\[3pt]
213 & Zenodo initialization without an analysis & 2023-08 & Needs human & API-design decision. \\[3pt]
206 & \code{environment.yml} for conda & 2023-08 & AI-tractable & Simple, self-contained. \\[3pt]
203 & E(B-V) contribution in likelihood.py & 2023-08 & Needs human & Physics-modeling feature. \\[3pt]
202 & dynesty/ultranest samplers hang indefinitely & 2023-08 & Needs human & Long-standing, hard sampler-hang bug. \\[3pt]
196 & Tests for \code{nmma/eos} & 2023-08 & AI-tractable & Matches this audit's \#1 coverage gap. \\[3pt]
188 & debug flag & 2023-08 & AI-tractable & Small, scoped CLI feature. \\[3pt]
184 & Release Notes for Version changes & 2023-07 & AI-tractable & Doc/process writing (once policy set). \\[3pt]
183 & Versioning schema & 2023-07 & Needs human & Policy decision. \\[3pt]
181 & Add compatibility list to documentation & 2023-07 & AI-tractable & Doc addition. \\[3pt]
159 & Move \code{best\_fit.json} into result JSON & 2023-07 & Stale (\code{wontfix}) & Already tagged; needs closing. \\[3pt]
142 & Plot multiple models onto fit lightcurve & 2023-06 & AI-tractable & \code{good first issue}, reference linked. \\[3pt]
139 & Strategy for adding models & 2023-06 & Needs human & Process/governance discussion. \\[3pt]
136 & Zenodo quickstart docs & 2023-06 & AI-tractable & Doc-writing task. \\[3pt]
125 & \code{svdmodel\_benchmark} multiprocessing bug & 2023-06 & Stale, triage & 3+ years old, unclear relevance. \\[3pt]
100 & Sequential Bayesian update for Hubble constant & 2023-06 & Needs human & Statistical-method design question. \\[3pt]
83 & Additional filter enforcement & 2023-05 & Stale, triage & Old, thin description. \\[3pt]
81 & Sample times double call & 2023-04 & Needs human & Abandoned PR \#82 attempt; needs maintainer decision on approach. \\[3pt]
78 & Improve \code{gwem\_resampling} & 2023-03 & Needs human & Vague scope, needs statistical design decision. \\[3pt]
20 & Sampler parameter tuning & 2022-04 & Stale, triage & Very old, likely superseded by \#202. \\
\bottomrule
\end{longtable}
\endgroup

\noindent\textbf{Issue summary:} 15 AI-tractable $\cdot$ 17 needs-human $\cdot$ 4
stale/needs-triage-first (including 1 already \code{wontfix}).

\section{GitHub PR triage (8 open, as of 2026-08-23)}

\begingroup
\small
\begin{longtable}{@{}p{1.2cm}p{3.3cm}p{2.6cm}p{2.6cm}p{4.4cm}@{}}
\toprule
\textbf{PR} & \textbf{Title} & \textbf{Size} & \textbf{Status} & \textbf{Recommendation} \\
\midrule
\endhead
\#433 & use richpool interface & tiny (+5/-5) & mergeable, review required &
  Touches untested \code{mpi\_setup.py}; author asks for help validating it works.
  \textbf{Needs human.} \\[4pt]
\#409 & Fix radec units & large (69 files, +2366/-1539) & conflicting, review
  required & Physics-critical correctness fix, wide blast radius, needs rebase.
  \textbf{Needs human.} \\[4pt]
\#397 & Generalizing Light Curve to Tensor Loading & medium (5 files, +596/-395) &
  conflicting, stale since Jul 2025 & Substantial refactor, needs design opinion
  and conflict resolution. \textbf{Needs human.} \\[4pt]
\#379 & LANL1D model addition & small (+37/-2) & conflicting, previously approved &
  New physics model; needs science review, has drifted into conflicts since
  approval. \textbf{Needs human} (not the pure fast win it once looked like). \\[4pt]
\#276 & add apptainer definition file & tiny (+25/-0) & mergeable, review required &
  Small, self-contained, already sanity-checked by author. \textbf{AI-tractable /
  fast win.} \\[4pt]
\#215 & multi\_model\_analysis command & draft, conflicting, small (+29/-0) &
  stale since Sep 2023 & Needs design review of interaction with existing
  multi-model support. \textbf{Needs human.} \\[4pt]
\#105 & Remove old Bu2019lm pickle / install updates & draft, conflicting,
  \textbf{+3.38M/-4.8k lines, 222 files} & stale since Jun 2023 & Binary/pickle
  files likely diffed as text; needs a maintainer to assess before review is even
  possible. \textbf{Needs human --- close or rebase from scratch.} \\[4pt]
\#82 & removed sample\_times arg in generate\_lightcurve & conflicting,
  +3.37M/-2.6k lines, 205 files & stale since Jun 2023 & Same red flag as \#105;
  not all call sites validated. Maps to issue \textbf{\#81} --- intent still
  valid, worth redoing small. \textbf{Needs human to decide close/redo; the redo
  itself is AI-tractable.} \\
\bottomrule
\end{longtable}
\endgroup

\noindent\textbf{PR summary:} 1 AI-tractable (fast win) $\cdot$ 7 needs-human. The
PR queue skews heavily toward ``needs human'' because the small/easy ones get
merged quickly --- what's left is either physics-critical, has merge conflicts, or
both.

\section{Summary counts}

\begin{center}
\small
\begin{tabular}{@{}lr@{}}
\toprule
\textbf{Category} & \textbf{Count} \\
\midrule
Headline findings & 11 \\
P0 / P1 / P2 action items & 7 / 7 / 6 \\
Needs human/expert judgment (\S3.1) & 17 \\
AI-agent-tractable (\S3.2) & $\sim$20 \\
Open issues --- AI-tractable & 15 \\
Open issues --- needs human & 17 \\
Open issues --- stale/needs triage & 4 \\
Open PRs --- AI-tractable & 1 \\
Open PRs --- needs human & 7 \\
\bottomrule
\end{tabular}
\end{center}

\noindent\textbf{Bottom line for the retreat:} the core likelihood abstraction
(\code{NMMALikelihoodMixin}/\code{MultiMessengerLikelihood}) is sound and doesn't
need architectural surgery. The real risk is concentrated in four places: (1) four
physics-critical/infra modules (\code{gw}, \code{eos}, \code{mlmodel},
\code{mpi\_setup.py}) with little-to-no test coverage, (2) two shipping-but-untested
bugs plus a $\sim$50--65-item FIXME backlog encoding real open physics questions,
(3) two actively-broken doc entry points that fail a new user on their very first
command, and (4) a backlog of $\sim$20 stale-to-old issues and several
multi-hundred-file PRs with merge conflicts that need a maintainer triage pass
before anyone --- human or AI --- can act on them productively. The AI-tractable
lists above are all good candidates to hand to coding agents \emph{during} the
retreat, freeing human time for the architecture/physics/triage decisions that
only the team can make.

\section{Suggested retreat schedule sketch}

\begin{itemize}
\item \textbf{Day 1 AM}: Triage session --- walk through \S5/\S6 as a group, assign
  owners, close what's dead (\#159, re-triage \#125/\#83/\#20), merge the fast win
  (\#276), decide what to do with \#379/\#409/\#397 given their merge-conflict state.
\item \textbf{Day 1 PM -- Day 2}: P0 items (\S2) --- coverage scaffolding for
  \code{eos}/\code{gw}/\code{mlmodel}/\code{population}/\code{mpi\_setup.py}, the
  two concrete bug fixes, bare-except and doc fixes, tooling/CI-pin bumps. Good
  pairing exercise: human defines what a correct test looks like, agent writes the
  boilerplate. Run \code{flake8 --max-complexity 15} for real on day 1 to
  corroborate \S4.2.
\item \textbf{Day 3}: Physics/design discussions from \S3.1 ---
  \code{conversion.py} FIXMEs, the \code{mlmodel/} maintain-vs-deprecate call,
  \code{core\_model\_name}, the pickle-dump versioning question, versioning/release
  policy. This is the block that most needs everyone in the room together.
\item \textbf{Day 4}: Work through P1 items and the agent-tractable issue list
  (\S3.2/\S5) in parallel --- pair each person with an agent session tackling one
  mechanical item at a time, human reviews before merge.
\item \textbf{Day 5}: Wrap-up --- decide the 1.1 release scope (\#429) now that the
  blockers above are visible, write release notes/changelog policy (\#183/\#184),
  retro on what worked.
\end{itemize}

\end{document}
